% !TeX root = ../papers/nc-p3.tex
\section{二重メビウスグラフと二次因子による有限鎖判定}
\label{sec:doubled-binary-pencils}
前二節と同じ三次曲線の残余で、埋込み次元が2以下とする。
前節の $E$ は隣接位置での $f$ の行列式を取ったスキームとする。
\begin{lemma}\label{lem:cubic-secant-isomorphism}
各辺の割線因子の対応 $M(d)g=0$ は、二つの平面三次スキームの同型のグラフである。
さらに $E$ の任意の成分・無限小構造を含め、$D_j(x)\mapsto S D_{j-2}(x)$ はこのグラフに入る。
\end{lemma}
\begin{proof}
三次式が非零であることは前節で証明した。
ある $d$ で $M(d)$ の核が二次元以上なら、$d S^1$ に支持される独立な二つの関係
$[z,t]d(z)g(t)$ ができ、定理\ref{thm:twisted-cubic-line-resolution}の二次での唯一性に反する。
逆向きも同じである。よって零点上では単元の二次小行列式による消去で射影がスキームとして同型になる。
$E$ 自身の隣の位置への移動も前節の単元による核の構成でスキーム同型である。
順序付き二点の全平面積で、四テンソルの二点縮約を先頭の括弧で割った二形式は
関係空間に値を取る。これは括弧が非零の開集合で成り立つ多項式恒等式だからである。
この恒等式を $E$ に制限し、階数1の行列を局所分解すれば所定の二因子を得る。
先頭の括弧が $E$ 上で零因子となる場合もこれでよい。
\end{proof}

\begin{theorem}\label{thm:doubled-binary-pencils}
$E$ が $(1,1)$ グラフの二重化なら、窓は標準Koszul左右AS正則である。
被約な二グラフの場合も、次のいずれかで同じ結論となる。
\begin{enumerate}
\item 一つの成分を全位置へ運んだとき、その上で各隣接座標の変換が恒等または対合となる。
\item 全ての位置で二つのグラフが互いの逆変換のグラフとなる。
\end{enumerate}
第一条件は四移動が成分を交換しても八つの連続位置で調べられる。
\end{theorem}
\begin{proof}
成分を一本取り、$h_j:\mathbb P^1\to C$ をその上の各位置の射影同型とする。
二重グラフの場合、$\phi=h_jh_{j-1}^{-1}$ が恒等でも対合でもなければ、
順序を忘れる射の下で一般点の近くで二重の二次曲線が現れる。
実際、グラフと逆グラフは異なり、順序を忘れる射は対角外でエタール、
被約グラフの像は次数2で一般点上で一対一である。
従ってこの二重グラフ上で消える平面多項式は二次曲線の一般点で少なくとも二重に消える。
補題の非零の三次式では不可能である。よって二重化は第一条件に入る。

第一条件で $a_j=h_{j-1}^{-1}h_j$ は恒等または対合である。
$x\mapsto D_j(x)$ は $a_j$ が恒等なら対角二次曲線への同型、
非自明なら直線への二重被覆で対合は $a_j$ である。
辺のグラフが同型なので各 $x$ について $D_j(x)$ と $D_{j-2}(x)$ は互いを決め、$a_j=a_{j-2}$ である。
二つの変換を $a,b$ と書く。これらに共通固定点 $x$ があれば、全 $h_j(x)$ は同じ $p$ で、
$D_j(x)=2p$ は位置によらない。共通固定点がなければ、両者は非自明で異なり、
$ab$ の固定点は相異なる二点である。一点なら両変換が $ab$ を逆に共役するのでそこを固定するからである。
また共通固定点の非存在により $a,b$ はともにこの二点を交換する。
実際、一方が二点を個別に固定すれば、積も固定するため他方も個別に固定する。
この二点集合を $Z$ とすると $h_j(Z)=h_0(Z)$ で、$x\in Z$ では $D_j(x)$ が常に同じ二次因子である。

第二条件では二つの異なるグラフは $\phi,\phi^{-1}$ で、$\phi$ は対合でない。
両者が交わる点の座標は $\phi^2$ の固定点であり、いずれも $\phi$ 自身の固定点である。
二成分の共通点を一つ選んで全位置へ運ぶと、各辺でその両座標が等しいので $D_j(x)=2p$ はやはり一定である。

以上のいずれでも、ある一定因子 $D$ について全 $j$ の辺に $D\mapsto SD$ がある。
四移動により各整数 $m$ について $S^mD\mapsto S^{m+1}D$ も全辺にある。
これで両方向の鉛筆列ができ、系\ref{cor:twisted-cubic-compatible-pencils}による。
\end{proof}

\begin{proposition}\label{prop:decide-compatible-secants}
この三次曲線の残余で、基因子の次数2の鉛筆による両方向の整合列が存在するかどうかは、
四窓の三次式、核同型および四移動に対する有限な閉集合演算で決定できる。
この試験に通る窓は全次数で正則である。
\end{proposition}
\begin{proof}
被約な射影平面の四つのコピー上で各窓の出発三次曲線を $B_i$ とする。
四辺目の終点を四移動で最初の位置に戻す。辺の同型を $f_i:B_i\to B'_i$ と書く。
$Z_i=B_i$ から始め、四位置を同時に
\[
 Z_i\ \longleftarrow\ Z_i\cap f_i^{-1}(Z_{i+1}\cap B'_i)
                       \cap f_{i-1}(Z_{i-1})
\]
で更新する（最後の項は位置 $i$ の平面に取る）。閉埋込みの像と逆像なので全て閉で、
根基イデアルの計算と等号判定で反復できる。閉集合の降鎖条件により有限回で安定する。
任意の両方向列は常に残る。逆に安定後の各点には残った前点・次点があり、
辺の同型によりこれを両方向へ延ばせる。四移動をほどけばもとの整合列である。
\end{proof}
